% Version 2022-09-2022
% update – 161114 by Ken Arroyo Ohori: made spacing closer to Word template throughout, put proper quotes everywhere, removed spacing that could cause labels to be wrong, added non-breaking and inter-sentence spacing where applicable, removed explicit newlines
% update – 010819 by Dennis Wittich: made spacing and font size closer to Word template, updated references and refernces style
% update – 042319 by Dennis Wittich: font size of captions set to 'small', first author names are shortened, hyphenation fixed
% update – 010620 by Dennis Wittich: Footnotes alignment set to left
% update - 151220 by Clement Mallet: Template adapted for single blind abstract submissions
% update - 060321 by Christian Heipke: Template refined for single blind abstract submissions
% update - 090921 by Christian Heipke: Template refined for single blind abstract submissions
% update - 200922 by Christian Heipke: general template update
% update - 080124 by Christian Heipke: general template update

\documentclass{isprs} % isprs class modified 23-04-2019 (Dennis Wittich)
\usepackage{subfigure}
\usepackage{setspace}
\usepackage{geometry} % added 27-02-2014 Markus Englich
\usepackage{epstopdf}
\usepackage[labelsep=period]{caption}  % added 14-04-2016 Markus Englich - Recommendation by Sebastian Brocks
\usepackage[spanish,es-tabla,es-nodecimaldot]{babel} 
\usepackage[hang]{footmisc}
\def\footnotemargin{1em} % added 08-01-2020 Dennis Wittich


\usepackage[authoryear]{natbib}
\def\bibhang{0pt}

\geometry{a4paper, top=25mm, left=20mm, right=20mm, bottom=25mm, headsep=10mm, footskip=12mm} % added 27-02-2014 Markus Englich
\usepackage{enumitem}
\usepackage{url}
\graphicspath{{figures/}}

%\usepackage{isprs}
%\usepackage[perpage,para,symbol*]{footmisc}

%\renewcommand*{\thefootnote}{\fnsymbol{footnote}}
\captionsetup{justification=centering,font=normal} % thanks to Niclas Borlin 05-05-2016
\captionsetup[figure]{font=small} % added 23-04-2019 Dennis Wittich
\captionsetup[table]{font=small} % added 23-04-2019 Dennis Wittich

\begin{document}

\title{Democratización de la Información Metropolitana: Apimetro como Infraestructura de Datos Abiertos y Motor de Análisis Topológico de Red en la ZMVM
}

% KAO: Remove extra spacing
\author{
 Hernán Galileo Cabrera Garibaldi\textsuperscript{1},
 David López\textsuperscript{2}
}

% KAO: Remove extra newline
\address{
	\textsuperscript{1 }Facultad de Estudios Superiores Acatlán, Universidad Nacional Autónoma de México, 53150 Naucalpan de Juárez, México\\
	313326626\@pcpuma.acatlan.unam.mx\\
	\textsuperscript{2 }Instituto de Ingeniería, Universidad Nacional Autónoma de México, 04510 Ciudad de México, México\\
	dlopezfl\@iingen.unam.mx
}

% If the corresponding author is NOT the final author, always add a % space before the subsequent comma, i.e.
% first author name\textsuperscript{a,}\thanks{Corresponding author} , % second author name \textsuperscript{b}, etc.
% thanks to Niclas Borlin 05-05-2016
% information on the corresponding author should not be used any longer and has been commented out
% C. Heipke, Jan 03,2024

% the use of the information of commissions and working groups should not be used any longer and has been commented out
% C. Heipke, Sept. 20,2022
%\commission{XX, }{YY} %This field is optional. If filled, XX and YY should be replaced by adequate numbers. See https://www2.isprs.org/commissions/
%\workinggroup{XX/YY} %This field is optional.
%\icwg{}   %This field is optional.


\abstract{
The Mexico City Metropolitan Area (ZMVM) operates one of the largest multimodal
transit networks in Latin America, yet access to unified geospatial data for this
network remains a persistent barrier: official GTFS feeds and OpenStreetMap data
are heterogeneous, institutionally fragmented, and technically demanding to use.
This paper presents Apimetro, a free and open-source infrastructure that normalizes
and integrates both sources into a unified REST service delivering GeoJSON, reducing
the technical barrier for urban researchers and practitioners.  Built on this
infrastructure, the \textit{VFTModel} analytical engine computes three network
efficiency indicators: spatial coverage, deviation factor, and capillary strength.
Results show that Milpa Alta borough records only 11.75\% pedestrian coverage and
that the average inter-borough trip covers 54\% more distance than the geometric
minimum.  The pattern is applicable in cities with available GTFS data and
OpenStreetMap coverage.
}

\keywords{Transporte Público, GTFS, OpenStreetMap, GeoJSON, Análisis topológico, ZMVM, Democratización de datos.}
\maketitle
\sloppy

\section{Introducción}
\label{introduccion}

La ZMVM opera una de las redes de transporte público más complejas del continente
bajo un patrón predominantemente radial que genera ineficiencias funcionales
documentadas \citep{padilla2022expansion, molinero2002}.  El estudio cuantitativo
de esta red enfrenta un problema recurrente de homologación: los archivos GTFS de
la Secretaría de Movilidad (SEMOVI) presentan severas discontinuidades ---rutas
inconclusas, atributos incompatibles entre esquemas y exclusión de sistemas
periféricos como el Mexibús \citep{semovi2020diagnostico}--- al cruzar las fronteras
entre la Ciudad de México y los municipios del Estado de México.  Aunque el derecho
de acceso a la información está constitucionalmente garantizado en México
\citep{cpeum_art6}, la apertura formal del dato no equivale a su democratización
\citep{gurstein2011open}: si su interpretación exige habilidades avanzadas de
ingeniería de software, el dato permanece inaccesible para la mayoría.

Ante esta fragmentación, OpenStreetMap (OSM) emerge como un recurso complementario:
su cartografía colaborativa produce datos de calidad comparable a los oficiales
\citep{haklay2010good, barrington2017world} y cubre los sistemas periféricos que los
GTFS gubernamentales omiten.  El presente trabajo expone cómo Apimetro, una
infraestructura de software libre, integra ambas fuentes ---GTFS oficiales y
OSM--- para construir un modelo de red metropolitana unificado expuesto como API
REST en formato GeoJSON, reduciendo la barrera técnica de entrada para
investigadores y profesionales urbanos.

El artículo se organiza como sigue: la sección~2 describe la arquitectura de
Apimetro y el motor analítico \textit{VFTModel}; la sección~3 presenta los
resultados sobre la red de la ZMVM; la sección~4 ofrece conclusiones e
implicaciones.


\section{Metodología}\label{sec:Metodologia}

El ecosistema Apimetro opera mediante una arquitectura desacoplada en dos capas.
La primera capa corresponde al servicio de unificación de datos espaciales,
construido en el lenguaje de programación Go ---que garantiza alta eficiencia en el
procesamiento concurrente de flujos masivos de datos \citep{goauthors2025}--- en combinación con una base
de datos relacional espacial PostgreSQL equipada con la extensión PostGIS.  Apimetro
ingiere los archivos GTFS oficiales de SEMOVI y los normaliza: homologa
campos faltantes y unifica formatos entre fuentes para que sean
directamente comparables.  Para resolver los vacíos de información del Estado
de México y corregir trazos geométricos gubernamentales discontinuos ---como el caso
del Trolebús Línea~13 o los sistemas de transporte metropolitano del Estado de
México---, Apimetro utiliza de forma estratégica la cartografía colaborativa de
OpenStreetMap como fuente de validación cruzada y datos complementarios.

Mediante consultas SQL espaciales, Apimetro unifica ambas fuentes y genera
una base de datos homogénea: un único formato donde geometrías,
identificadores y atributos son consistentes entre sistemas.  La segunda capa de Apimetro corresponde al motor
analítico \textit{VFTModel}, que consume los servicios REST resultantes para
construir grafos matemáticos y calcular indicadores de eficiencia de red.

La contribución de OSM al \textit{pipeline} de extracción y tratamiento de datos
de Apimetro opera en tres niveles concretos:

(1)~\textbf{geometrías de ruta}, cuando los trazos GTFS presentan discontinuidades
o resolución insuficiente, Apimetro consulta las relaciones OSM de tipo
\texttt{route=bus}, \texttt{route=subway} y \texttt{route=light\_rail} para obtener
geometrías continuas de alta resolución; (2)~\textbf{cobertura periférica}, los
municipios conurbados del Estado de México carecen en gran medida de archivos GTFS
publicados, por lo que la comunidad OSM cubre los recorridos que las fuentes
oficiales omiten \citep{haklay2010good} ---las rutas del Mexibús y los autobuses
periféricos se obtienen íntegramente de OSM; y (3)~\textbf{validación cruzada},
Apimetro emplea la distancia de Hausdorff para detectar inconsistencias geométricas
entre trazas GTFS y OSM \citep{hausdorff1914, neis2012street}, marcando para
revisión las discrepancias superiores a un umbral configurable.

El sistema expone la red multimodal completa ---Metro, Metrobús, Cablebús, RTP,
Trolebús, Tren Ligero, Tren Suburbano, Mexibús, Colectivo, Microbús y
Autobús--- dividida en dos entidades relacionales: \textit{estaciones} (puntos con
jerarquía de transporte y nodos CETRAM identificados) y \textit{líneas} (geometrías
\texttt{MultiLineString} segmentadas tramo a tramo entre estaciones).  Un proceso de \textit{snapping} con umbral de $\leq 50$~m resuelve las
interconexiones entre sistemas por proximidad geográfica.  De los 11~sistemas integrados, el 88\% corresponde
a transporte de superficie y el 12\% a infraestructura de vía exclusiva.  La salida
se expone públicamente a través de servicios REST en formato GeoJSON nativo
desde \texttt{apimetro.dev}.

Para demostrar la capacidad analítica de la infraestructura, se desarrolló el
\textit{VFTModel} (\textit{Vanishing Fig-Tree Model})%
\footnote{Repositorio: \url{https://github.com/galigaribaldi/VFTModel};
documentación: \url{https://galigaribaldi.github.io/VFTModel/notebooks/intro.html}},
un motor analítico en Python que consume directamente los servicios GeoJSON de
Apimetro.  Utilizando librerías de análisis espacial y redes complejas (NetworkX,
GeoPandas, Momepy), el modelo transforma las geometrías vectoriales en grafos
matemáticos ruteables y calcula tres indicadores de eficiencia de red
\citep{fortin2016innovative}: (1)~\textbf{cobertura espacial},
porcentaje de superficie urbana dentro del radio peatonal de 800~m de cada estación;
(2)~\textbf{factor de desviación}, exceso de la distancia real en red sobre la
distancia euclidiana mínima entre dos puntos; y (3)~\textbf{fuerza capilar},
concentración de flujos en nodos del grafo con escasas rutas alternativas.  Los
resultados de su aplicación a la ZMVM se detallan en la
sección~\ref{sec:Resultados}; los datos se exportan como capas vectoriales
GeoPackage (\texttt{.gpkg}) directamente visualizables en QGIS.

\section{Resultados}\label{sec:Resultados}

Los tres indicadores calculados por el VFTModel sobre la red de Apimetro arrojan
los siguientes resultados para la ZMVM:
(1)~\textbf{cobertura espacial ($C$)}, isócronas peatonales de 800~m alrededor de
cada estación ---Milpa Alta presenta solo el 11.75\% de su superficie dentro del
radio de cobertura (véase Anexo);
(2)~\textbf{factor de desviación ($DF$)}, cociente entre la
distancia real en red y la distancia euclidiana mínima ---el promedio en
100~rutas interalcaldías arroja $DF=1.545$, es decir, cada viaje recorre 54\%
más de lo necesario (caso extremo San Pedro Atocapan, Alcaldía Milpa Alta, a
Santa Catarina Yecahuizotl, Alcaldía Tláhuac: $DF=2.52$; véase Anexo); y
(3)~\textbf{fuerza capilar ($FC_H$)}, índice que cuantifica la concentración de
flujos en nodos del grafo con pocas rutas alternativas ---el nodo máximo registra
$FC_H=130$ en la zona de Hospitales Sur (véase Anexo).

En total, el motor procesa 22\,766 entidades espaciales que se consolidan en un
grafo de 11\,115~nodos interconectados por 45\,679~aristas: 25\,197 de servicio
real, distribuidas en 11~sistemas (RTP:~11\,041; CC:~10\,603; TROLE:~1\,506;
MB:~792; METRO:~528; MEXIBÚS:~391; PUMABUS:~234; TL:~52; CBB:~30;
MEXICABLE:~17; SUB:~3), y 20\,482 de transbordo peatonal, todo en un volumen
de 38~MB de datos procesados.  Estos resultados ilustran cómo Apimetro facilita
el acceso a datos de red normalizados para cualquier usuario, sin necesidad de
programar, limpiar archivos de texto ni reparar trazos geométricos rotos.
Los recursos interactivos completos ---red unificada y panel de consulta REST--- están
disponibles en el Anexo.

\section{Conclusiones}
\label{sec:Conclusiones}

La combinación de datos gubernamentales (GTFS) y cartografía colaborativa
(OpenStreetMap) que articula Apimetro ofrece una vía concreta para reducir las
barreras técnicas en la investigación de redes de transporte urbano.  Al exponer
una base PostGIS compleja como API REST en formato GeoJSON, cualquier urbanista,
geógrafo o desarrollador puede integrar la red multimodal de la ZMVM en herramientas
como QGIS, Kepler.gl o Leaflet sin necesidad de programar ni limpiar datos.

El patrón arquitectónico de Apimetro ---ingesta de GTFS, normalización, unificación
con OSM y exposición vía API REST--- es reproducible en cualquier ciudad que disponga
de archivos GTFS y cobertura de OpenStreetMap, lo que lo convierte en un recurso
potencialmente útil para otras zonas metropolitanas con problemáticas similares de
fragmentación de datos.

Se invita a la comunidad de OpenStreetMap en Latinoamérica a colaborar en tres
frentes: (1)~\textbf{mapatones temáticos}, sesiones de mapeo enfocadas en sistemas
de transporte periférico cuya cobertura OSM es aún incompleta en municipios
conurbados como Ecatepec, Nezahualcóyotl y Chimalhuacán; (2)~\textbf{validación
bidireccional}, publicar como datos abiertos los resultados del proceso de
\textit{conflation} de Apimetro para que la comunidad corrija errores detectados
tanto en OSM como en los GTFS oficiales; y (3)~\textbf{replicabilidad regional},
documentar el patrón GTFS+OSM\,$\to$\,API REST como modelo de referencia para
otras zonas metropolitanas latinoamericanas (p.e.\ Bogotá, São Paulo, Buenos Aires)
que enfrenten fragmentación de datos de transporte.

La experiencia de Apimetro sugiere que los datos gubernamentales abiertos y la
cartografía colaborativa son esfuerzos complementarios: la calidad de uno potencia
la utilidad del otro.  Entre las líneas de trabajo futuro se identifican la
incorporación de indicadores personalizables por el usuario, la mejora de la
integración OSM--GTFS en municipios con baja cobertura colaborativa, y la
documentación del sistema para públicos no técnicos.

{
	\begin{spacing}{1.17}
		\normalsize
		\bibliography{ISPRSguidelines_authors}
	\end{spacing}
}

\section*{Anexo}

\begin{itemize}
  \item \textbf{Dashboard Apimetro} (red multimodal ZMVM y panel de consulta REST):\\
        \url{https://public.tableau.com/views/DashboardRedApimetro2025/DashboardResultadosApimetro}

  \item \textbf{VFTModel --- Factor de desviación} (distribución en 100 rutas interalcaldías):\\
        \url{https://public.tableau.com/views/RedVFTModelZMVM2025-2027FacotrdeDeviacin/PlantillaDashboardUnigrafico}

  \item \textbf{VFTModel --- Fuerza capilar} (top estaciones por concentración de flujos):\\
        \url{https://public.tableau.com/views/RedVFTModelZMVM2025-2027FuerzaCapilarTopEstacionesTabla/PlantillaDashboardUnigrafico2}

  \item \textbf{Mapa de calor --- Cobertura espacial peatonal} (isócrona 800~m, VFTModel):\\
        \url{https://github.com/galigaribaldi/Transport-gis-zmvm-mjg/releases/tag/v2026-2-col-mapas-cobertura-calor}
\end{itemize}

\end{document}
